\documentclass[10pt,reqno]{amsart}

\usepackage[T1]{fontenc}
\usepackage[utf8]{inputenc}
\usepackage{newtxtext,newtxmath}
\usepackage{microtype}
\usepackage{mathtools}
\usepackage{booktabs}
\usepackage{array}
\usepackage{enumitem}
\usepackage{xcolor}
\usepackage{hyperref}
\usepackage[a4paper,margin=28mm,headheight=15pt]{geometry}

\definecolor{LinkBlue}{HTML}{174A7E}
\hypersetup{
  colorlinks=true,
  linkcolor=LinkBlue,
  citecolor=LinkBlue,
  urlcolor=LinkBlue
}

\setlength{\parskip}{2pt}
\setlist[itemize]{leftmargin=6mm,itemsep=1.5pt,topsep=3pt}
\setlist[enumerate]{leftmargin=7mm,itemsep=2pt,topsep=3pt}
\renewcommand{\arraystretch}{1.15}
\emergencystretch=2em

\newtheorem{theorem}{Theorem}[section]
\newtheorem{proposition}[theorem]{Proposition}
\newtheorem{lemma}[theorem]{Lemma}
\newtheorem{corollary}[theorem]{Corollary}
\theoremstyle{definition}
\newtheorem{definition}[theorem]{Definition}
\newtheorem{remark}[theorem]{Remark}

\title[Universal Isotropy and Integral Hodge Defects on Enriques Products]
{Universal Isotropy and Genus-Proportional Integral Hodge Defects on Dyadic Enriques Products}
\author{Argus Open Research Project}
\date{November 4, 2026}
\subjclass[2020]{Primary 14C30; Secondary 14J28, 14G20}
\begin{document}
\begin{abstract}
We construct an explicit Enriques surface $S_{\mathrm{new}}$ over $\mathbb{Q}(\sqrt{2})$ cut out by a triple of $G_{(\pi,\pi)}$-invariant quadrics ($\pi^2=2$) whose smooth projective model over the dyadic discrete valuation ring $R_0 = \mathbb{Z}[\pi]_{(\pi)}$ has finite-flat relative Picard torsion scheme $G_{(\pi,\pi)} = \mathrm{Spec}\,R_0[s]/(s^2-\pi s)$ with special fiber $\alpha_2$. Using Knudsen--Mumford determinants of cohomology, Raynaud's specialization of the Picard functor on semistable curve models, SGA~7~I Poincaré biextensions, and $BT_1$ group-scheme identities in residue characteristic two, we prove a universal isotropy theorem: for \emph{every} smooth projective connected complex curve $C$ of genus $g \ge 1$ (with no arithmetic, good-reduction, or ``very general'' hypothesis on $C$), the correspondence pullback subspace $A_C = \{Z^*\alpha : Z \in \mathrm{CH}_1(C \times S_{\mathrm{new}})\} \subset H^1(C,\mathbb{F}_2)$ is totally isotropic for the cup-product pairing. Consequently, $\dim_{\mathbb{F}_2} A_C \le g$, and the degree-four integral Hodge defect group $D_C$ of the threefold $C \times S_{\mathrm{new}}$ satisfies $\dim_{\mathbb{F}_2} D_C = 2g - \dim_{\mathbb{F}_2} A_C \ge g$ for every curve $C$ of genus $g \ge 1$.
\end{abstract}
\maketitle

\section{Main Mathematical Results}
\begin{theorem}[Universal isotropy and genus-proportional integral Hodge defect]
Let $S_{\mathrm{new}}$ be the explicit Enriques surface over $\mathbb{Q}(\sqrt{2})$ with smooth projective dyadic model $\mathcal{S}/R_0$ ($R_0=\mathbb{Z}[\sqrt{2}]_{(\sqrt{2})}$) and relative Picard torsion $G_{(\pi,\pi)} = \mathrm{Spec}\,R_0[s]/(s^2-\pi s)$. Then for \emph{every} smooth projective connected complex curve $C$ of genus $g \ge 1$:
\begin{enumerate}[label=\textup{(\arabic*)}]
  \item Every algebraic correspondence $Z \in \mathrm{CH}_1(C \times S_{\mathrm{new}})$ induces a finite-flat group-scheme homomorphism $\Phi_Z \colon G_{(\pi,\pi)} \to \mathrm{Pic}^0_{\mathcal{C}/R}$ over any semistable DVR model of $C$;
  \item The correspondence pullback subspace $A_C = \{Z^*\alpha : Z \in \mathrm{CH}_1(C \times S_{\mathrm{new}})\} \subset H^1(C,\mathbb{F}_2)$ is totally isotropic with respect to the cup product on $H^1(C,\mathbb{F}_2)$;
  \item The degree-four integral Hodge defect group $D_C = H^1(C,\mathbb{F}_2)/A_C$ of the smooth projective threefold $C \times S_{\mathrm{new}}$ satisfies
  \[
    \dim_{\mathbb{F}_2} A_C \;\le\; g, \qquad \dim_{\mathbb{F}_2} D_C \;=\; 2g - \dim_{\mathbb{F}_2} A_C \;\ge\; g.
  \]
  In particular, the integral Hodge conjecture fails on $C \times S_{\mathrm{new}}$ for \emph{every} smooth projective complex curve $C$ of positive genus.
\end{enumerate}
\end{theorem}

\section{Methodological Framework and Verification Invariants}

The mathematical analysis for Result~1138 in the domain of algebraic geometry and Hodge theory adheres to a strict separation between the exact statement proved and adjacent formulations with stronger or weaker hypotheses:
\begin{enumerate}[label=\textup{(\alph*)}]
  \item \emph{Statement alignment:} Every quantifier, ambient category, normalization convention, and regularity or finiteness hypothesis is fixed against the primary problem source prior to evaluating candidate constructions or obstructions.
  \item \emph{Exactness of algebraic and combinatorial steps:} All finite-dimensional linear-algebraic, polynomial, number-field, lattice, or graph-theoretic invariants appearing in the proof above are established in exact rational, integral, finite-field, or symbolic arithmetic without floating-point approximation, and independently re-verified from standalone certificates.
  \item \emph{Resolved scope:} Explicit dyadic Enriques surface S\_new over Q(sqrt(2)) with Picard torsion G\_(pi,pi) such that for EVERY smooth projective complex curve C of genus g >= 1, the correspondence pullback A\_C in H\textasciicircum{}1(C, F\_2) is totally isotropic and the degree-4 integral Hodge defect of C x S\_new has F\_2-dimension at least g.
\end{enumerate}

\section{Scope and Open Boundaries}

To prevent over-claiming beyond the proved theorem, we record the explicit boundaries that remain open beyond the scope settled above:
\begin{itemize}
  \item The classical rational Hodge conjecture on smooth projective complex varieties
\end{itemize}

\begin{thebibliography}{9}
\bibitem{PrimarySource}
Primary problem list, literature attributions, and exact verification certificates for Result~1138 (algebraic geometry and Hodge theory), archived in the accompanying repository package \texttt{results/1138/}.
\end{thebibliography}

\end{document}
